%Document Preamble of Monabbati groups for Ferdowsi univercity of mashhad
%-------------------------------------------------------------------------------------------------------------------
%  دستوراتی برای اصلاح تولید نمایه
\immediate\write18{xindy -L persian -C utf8 -M texindy \jobname.idx}
%  تعریف کلاس سند
%  Kashida: برای تنظیم کشیدگی خودکار متن(این بسته باعث اختلال در خط نستعلیق می‌شود، openany: برای عدم شروع از صفحات فرد، ar: برای تنظیم پایان‌نامه، phd: برای تنظیم رساله، fleqn: برای؟؟؟
\documentclass[a4paper, 11pt, openany, oneside, arshad, Kashida]{FThesis}
%, Kashida
\usepackage[font={footnotesize }]{caption}
% در ورژن جدید زی‌پرشین برای تایپ متن‌های ریاضی، این سه بسته، حتماً باید فراخوانی شود
%\usepackage{amsthm, amssymb, amsmath}
% بسته‌ای برای تنطیم حاشیه‌های بالا، پایین، چپ و راست صفحه

\usepackage[top=20mm, bottom=20mm, left=25mm, right=35mm]{geometry}
% بسته‌‌ای برای ظاهر شدن شکل‌ها و تصاویر متن
\usepackage{graphicx}
% بسته‌ای برای رسم کادر\usepackage{graphicx,subfigure}
\usepackage{graphicx,subfigure}
%تداخل بسته subfigure با tocloft
\PassOptionsToPackage{subfigure}{tocloft}
\usepackage{fixmetodonotes}
%\usepackage{float}

\usepackage{amsthm}
\usepackage{amsmath}

\usepackage{indentfirst}
\usepackage{fancyhdr}
\usepackage{amsthm, amssymb, amsmath}
\usepackage{framed} 
% بسته‌‌ای برای چاپ شدن خودکار تعداد صفحات در صفحه «معرفی پایان‌نامه»
\usepackage{lastpage}
% بسته‌‌ای برای ایجاد دیاگرام‌های مختلف
\usepackage[all]{xy}
% بسته‌ و دستوراتی برای ایجاد لینک‌های رنگی با امکان جهش
\usepackage[pagebackref=false, colorlinks, linkcolor=black, citecolor=magenta]{hyperref}
% چنانچه قصد پرینت گرفتن نوشته خود را دارید، خط بالا را غیرفعال و  از دستور زیر استفاده کنید چون در صورت استفاده از دستور زیر‌‌، 
% لینک‌ها به رنگ سیاه ظاهر خواهند شد که برای پرینت گرفتن، مناسب‌تر است
%\usepackage[pagebackref=false]{hyperref}
% بسته‌ لازم برای تنظیم سربرگ‌ها
\usepackage{fancyhdr}
% بسته‌ای برای ظاهر شدن «مراجع» و «نمایه» در فهرست مطالب
%‪ \usepackage[nottoc,notlof,notlot]{tocbibind}
% دستورات مربوط به ایجاد نمایه
\usepackage{makeidx}
\makeindex
% بسته‌هایی برای رسم شکل
%\usepackage{subfigure}


\usepackage[usenames,dvipsnames]{pstricks}
\usepackage{epsfig}
\usepackage{pst-grad} % For gradients
\usepackage{pst-plot} % For axes
\usepackage{pst-node}
%\usepackage{tikz}
%\usetikzlibrary{decorations.pathmorphing} % noisy shapes
%\usetikzlibrary{fit}					% fitting shapes to coordinates
%\usetikzlibrary{backgrounds}	% drawing the background after the foreground
% بسته‌ای که باعث ریست شدن شماره پانوشت‌ها در هر صفحه می‌شود

\usepackage{cite}
% بسته‌ای برای چندستونی نوشتن
\usepackage{multicol}
% بسته‌هایی برای الگوریتم
\usepackage{algorithm}
%\usepackage{algpseudocode}
%\usepackage{enumerate}
\usepackage{bidiftnxtra}
% بسته‌ای برای اضافه کردن کدهای متلب
%\usepackage[framed,numbered,autolinebreaks,useliterate]{mcode}

% بسته‌ای برای رسم دیاگرام
\usepackage[all]{xy}

% بسته‌هایی برای رسم گراف
\usepackage{tkz-berge}
\usepackage{tkz-graph}
% بسته‌ای برای اضافه کردن نمایه
\usepackage[xindy]{glossaries}
%\makeglossaries
%\glossarystyle{mylist}
%\def\glossaryname{واژه نامه}
%\newcommand{\mygls}[2]{\newglossaryentry{#1}{name=#1,description={#2}}}

% بسته‌ای برای حذف بخشی از سند
\usepackage{comment}


% بسته‌ای برای رسم فلوچارت

\usepackage{mathrsfs, setspace,  fancybox, subfigure,  hyperref}
%, breqn
% بسته breqn با بسته xy مشکل داره باید بسته pb-diagram رو چک کنم.
%%%%%%%%%%%%%%%%%%%%%%%%%%
% فراخوانی بسته زی‌پرشین و تعریف قلم فارسی و انگلیسی


\usepackage[Kashida]{xepersian}
\usepackage{xepersian}


\usepackage{tikz}
\usetikzlibrary{shapes,arrows}

\usepackage[noend]{algpseudocode}
%\usepackage{perpage}
%\MakePerPage{footnote}
\settextfont[Scale=1.2]{Nazanin}%Persian Modern
% از revision 118 زی‌پرشین به بعد، وارد کردن دستور زیر لازم نیست. توجه داشته باشید که در صورت  غیرفعال کردن این دستور،
% از فونت پیش‌فرض لاتک برای کلمات انگلیسی استفاده خواهد شد.
%\setlatintextfont[ExternalLocation,BoldFont={lmroman10-bold},BoldItalicFont={lmroman10-bolditalic},ItalicFont={lmroman10-italic}]{lmroman10-regular}
%%%%%%%%%%%%%%%%%%%%%%%%%%
% چنانچه می‌خواهید اعداد در فرمول‌ها، انگلیسی باشد، خط زیر را غیرفعال کنید
%\setdigitfont[Scale=1.1]{Parsi Digits}
%%%%%%%%%%%%%%%%%%%%%%%%%%

% تعریف قلم‌های فارسی و انگلیسی اضافی برای استفاده در بعضی از قسمت‌های متن
\defpersianfont\nastaliq[Scale=2]{IranNastaliq}
\defpersianfont\nastaliqt[Scale=4]{IranNastaliq}
\defpersianfont\sols[Scale=1.5]{XB Sols}
\defpersianfont\chapternumber[Scale=3]{XB Niloofar}
\defpersianfont\titr[Scale=1]{XB Titre}
\setlatintextfont{Times New Roman}
\defpersianfont\nazanint[Scale=1]{Nazanin}
%%%%%%%%%%%%%%%%%%%%%%%%%%
% دستوری برای حذف کلمه «چکیده»
%\renewcommand{\abstractname}{}
% دستوری برای حذف کلمه «abstract»
%\renewcommand{\latinabstract}{}
% دستوری برای تغییر نام کلمه «اثبات» به «برهان»
\renewcommand\proofname{\textbf{برهان}}
% دستوری برای تغییر نام کلمه «کتاب‌نامه» به «مراجع»
\renewcommand{\bibname}{مراجع}
% دستوری برای تعریف واژه‌نامه انگلیسی به فارسی
\newcommand\persiangloss[2]{#1\dotfill\lr{#2}\\}
% دستوری برای تعریف واژه‌نامه فارسی به انگلیسی 
\newcommand\englishgloss[2]{#2\dotfill\lr{#1}\\}
%%%%%%%%%%%%%%%%%%%%%%%%%%
% تعریف و نحوه ظاهر شدن عنوان قضیه‌ها، تعریف‌ها، مثال‌ها و ...


\tikzstyle{vertex}=[circle, draw, inner sep=0pt, minimum size=15pt]
\usetikzlibrary{backgrounds,fit,decorations.pathreplacing}  % TikZ libraries
\usetikzlibrary{shapes,arrows}
\renewcommand{\labelitemi}{$\bullet$}
\newcommand{\vertex}{\node[vertex]}
\tikzstyle{vertexs}=[draw,minimum width=.5cm,minimum height=.5cm]
\usetikzlibrary{backgrounds,fit,decorations.pathreplacing,calc}
\newcommand{\ket}[1]{\ensuremath{\left|#1\right\rangle}} % Dirac Kets
\newcommand{\vertexs}{\node[vertexs]}
\theoremstyle{definition}
\newtheorem{definition}{تعریف}[section]
\newtheorem{example}[definition]{مثال}
\theoremstyle{theorem}
\newtheorem{theorem}[definition]{قضیه}
\newtheorem{lemma}[definition]{لم}
\newtheorem{proposition}[definition]{گزاره}
\newtheorem{corollary}[definition]{نتیجه}
\newtheorem{remark}[definition]{ملاحظه}
\newtheorem{treaty}[definition]{قرارداد}
\newtheorem{rema}[definition]{تبصره}
\newtheorem{notice}[definition]{توجه}
\newtheorem{hint}[definition]{تذکر}
\newtheorem{problem}[definition]{مساله}
\newtheorem{rem}[definition]{یادآوری}
\newtheorem{x}[definition]{}
\newenvironment{solution}{{\bf \noindent \\ حل. } \rm } 
% محیط‌های آزاد با و بدون شماره

\newcommand{\nameB}{}
\newtheorem*{envB}{\nameB}
\newenvironment{fenv}[1][پیش‌فرض]
{\renewcommand{\nameB}{#1}\begin{envB}}{\end{envB}}

\newcommand{\nameA}{}
\newtheorem{envA}{\nameA}
\newenvironment{fcenv}[1][پیش‌فرض]
{\renewcommand{\nameA}{#1}\begin{envA}}{\end{envA}}

%%%%%%%%%%%%%%%%%%%%%%%%%%
% تعریف دستورات جدید برای خلاصه نویسی و راحتی کا در هنگام تایپ فرمول‌های ریاضی
%\newcommand{\op}[1]{\mathop{#1}\limits}
%%%%%%%%%%%%%%%%%%%%%%%%%%%%
% دستورهایی برای سفارشی کردن سربرگ صفحات
\SepMark{-}


\csname@twosidetrue\endcsname
\pagestyle{fancy}
\fancyhf{} 
\fancyhead[OL,EL]{\thepage}
\fancyhead[OR,ER]{\small\leftmark}
%\fancyhead[ER]{\small\leftmark}
\renewcommand{\chaptermark}[1]{%
\markboth{#1}{}}
%\markboth{\thechapter.\ #1}{}}
%%%%%%%%%%%%%%%%%%%%%%%%%%%%%
% دستورهایی برای سفارشی کردن صفحات اول فصل‌ها
%\makeatletter
%\newcommand\mycustomraggedright{%
% \if@RTL\raggedleft%
% \else\raggedright%
% \fi}
%\def\@part[#1]#2{%
%\ifnum \c@secnumdepth >-2\relax
%\refstepcounter{part}%
%\addcontentsline{toc}{part}{\thepart\hspace{1em}#1}%
%\else
%\addcontentsline{toc}{part}{#1}%
%\fi
%\markboth{}{}%
%{\centering
%\interlinepenalty \@M
%\ifnum \c@secnumdepth >-2\relax
% \huge\bfseries \partname\nobreakspace\thepart
%\par
%\vskip 40\p@
%\fi
%\Huge\bfseries #2\par}%
%\@endpart}
%\def\@makechapterhead#1{%
%\vspace*{80\p@}%
%{\parindent \z@ \mycustomraggedright %\@mycustomfont
%\ifnum \c@secnumdepth >\m@ne
%\if@mainmatter
% این دو جفت محیط center باعث وسط چین شدن عنوان فصل می‌شوند.
%\begin{center}
%\huge\bfseries\space {\chapternumber\thechapter}
%\end{center}
%\par\nobreak
%\vskip 20\p@
%\fi
%\fi
%\interlinepenalty\@M 
%\begin{center}
%\Huge \bfseries #1\par\nobreak
%\end{center}
%\vskip 50\p@
%}}
\makeatother
% دستوراتی برای تنظیم شماره لاتین برای پانوشت‌ها
\makeatletter
\def\@makeLTRfnmark{\hbox{\@textsuperscript{\latinfont\@thefnmark}}}
\renewcommand\@makefntext[1]{%
    \parindent 1em%
    \noindent
    \hb@xt@1.8em{\hss\if@RTL\@makefnmark\else\@makeLTRfnmark\fi}#1}
\makeatother
% دستری برای تنظیم فرمول‌ها در حالت نمایشی
\everymath\expandafter{\the\everymath \displaystyle}